\section{Proposed Framework}
\label{sec:framework}

The framework introduced in Section~\ref{sec:architecture} is evaluated along
four protocols that together characterise its operating envelope. Multi-class
attack-family classification (\S\ref{subsec:multiclass}) measures closed-set
accuracy. Zero-day detection via Suspicious Multi-Class Confidence
(\S\ref{subsec:zeroday}) measures graceful degradation on out-of-distribution
traffic. The anonymization ablation (\S\ref{subsec:anon}) measures the
privacy-utility tradeoff under feature-name obfuscation. The adversarial
robustness protocol (\S\ref{subsec:kerckhoffs}) measures evasion cost under a
Kerckhoffs-style threat model. All experiments use
\texttt{claude-haiku-4-5-20251001} as the rule generator and fix the global
hyperparameters $k_{\text{rules}}=5$, $n_{\text{top}}=5$, $n_{\text{seeds}}=3$,
and $n_{\text{rounds}}=5$ with seeds $\{42, 123, 456\}$.

\subsection{Multi-Class Classification Protocol}
\label{subsec:multiclass}

Multi-class classification is realised through three coupled components. Each
dataset is labelled with its full taxonomy of attack families. Labels are
stored alongside feature vectors in the Chroma vector store so that retrieved
neighbours carry class provenance. The classification prompt lists the
complete set of $K$ known attack classes and instructs Claude Haiku 4.5 to
assign one of these labels. Majority voting across the top-$k$ retrieved
neighbours' labels supplies a soft prior to the model as in-context evidence.
Following classification, the rule-generation module synthesises one
decision-rule vector per detected class. Each vector consists of
$k_{\text{rules}}=5$ threshold conditions over the features ranked highest by
permutation importance ($n_{\text{top}}=5$), emitted as JSON tuples
\texttt{(feature\_name, value, op)} for downstream deployment as ACL-style
firewall entries.

\subsection{Zero-Day Detection Protocol}
\label{subsec:zeroday}

Rule-based detectors are closed-set by construction and cannot name a truly
novel attack class. This protocol assesses whether retrieval confidence
degrades gracefully when out-of-distribution traffic appears, enabling an
anomaly flag without a matching rule. For each dataset and a designated
withheld class $c^{*}$, all flows of $c^{*}$ are removed from the training
split and from the Chroma store. The remaining $K-1$ classes form the
closed-set training regime, and $c^{*}$ flows are reintroduced at inference
time alongside balanced known-class samples. Rather than asking the model to
label unseen traffic, we monitor the maximum softmax confidence of the
retrieved neighbour distribution. For a query $\mathbf{x}$, let $p_j$ denote
the proportion of retrieved neighbours labelled with class $j$. The SMCC
score is
\[
  \mathrm{SMCC}(\mathbf{x}) = \max_{j \in \mathcal{K}} p_j ,
\]
where $\mathcal{K}$ is the set of known classes. When
$\mathrm{SMCC}(\mathbf{x}) < \tau$, the query is flagged as a candidate
zero-day event. The threshold $\tau$ is swept from $0$ to $1$ and the
Zero-Day Detection Rate (ZDR) and False Positive Rate are recorded at each
point, producing the curves reported in \S\ref{subsec:threshold}. ZDR is the
fraction of withheld-class samples correctly flagged,
\[
  \mathrm{ZDR} = \frac{\mathrm{TP}_{c^{*}}}{\left| \mathcal{D}_{c^{*}} \right|} ,
\]
where $\mathcal{D}_{c^{*}}$ is the set of test flows belonging to $c^{*}$.

\subsection{Anonymization Ablation}
\label{subsec:anon}

The prompt necessarily exposes raw feature names such as
\texttt{Header\_Length} or \texttt{SrcAddr}, which can leak deployment topology
or vendor information. This ablation tests whether replacing semantic names
with opaque identifiers degrades detection. For each dataset, a bijective map
$\sigma: \text{feature\_name} \to f_i$ is constructed once and held fixed
across seeds. The model receives only the anonymized identifiers and the
numeric values. Two conditions are compared across $n_{\text{seeds}}=3$ random
seeds and $n_{\text{rounds}}=5$ refinement rounds. The \emph{Original}
condition uses semantic names. The \emph{Anonymized} condition uses opaque
identifiers. Performance is reported as macro-averaged F1 on a held-out test
set, and the utility cost of anonymization is
$\Delta F1 = \overline{F1}_{\text{anon}} - \overline{F1}_{\text{orig}}$.
Feature-set consistency across seeds is quantified by the Jaccard index over
the selected $k_{\text{rules}}=5$ features,
\[
  J = \frac{\left| \mathcal{F}_{\text{orig}} \cap \mathcal{F}_{\text{anon}} \right|}
           {\left| \mathcal{F}_{\text{orig}} \cup \mathcal{F}_{\text{anon}} \right|} .
\]

\subsection{Adversarial Robustness Protocol}
\label{subsec:kerckhoffs}

The robustness protocol assumes the Kerckhoffs principle. The attacker has
full knowledge of the generated rule set and applies the minimum feature
perturbation needed to evade detection. The perturbation budget $\varepsilon$
is expressed as a fraction of each feature's interquartile range, and the
adversary perturbs withheld-class flows toward the nearest known-class
centroid. The Exploitability Sweep Rate measures the fraction of
originally-flagged flows that evade detection at budget $\varepsilon$,
\[
  \mathrm{ESR}(\varepsilon) = \frac{\text{flows that evade at budget } \varepsilon}
                                   {\text{flows flagged at } \varepsilon = 0} .
\]
The summary metric is the critical budget $\varepsilon_{0.5}$ at which
$\mathrm{ESR}=0.5$, which quantifies the adversarial hardness of each
detector. Five conditions are compared. \texttt{LLM\_orig} uses semantic
feature names. \texttt{LLM\_anon} uses opaque identifiers. \texttt{DT\_d3}
and \texttt{DT\_d5} are Decision Trees at depths three and five.
\texttt{RF\_d5} is a depth-five Random Forest. Each condition is run for
three seeds, and the mean and standard deviation of $\varepsilon_{0.5}$ are
reported in Table~\ref{tab:kerckhoffs}.
